\section{Additional Tables}
\label{app:additional-tables}

This appendix reports the minimum-wage distribution tables, alternative-reference tables, comparison specifications, and income-scenario checks referenced in the main text. Table~\ref{tab:minimum-wage-bunching-2025} complements Figure~\ref{fig:minwage-distribution-formality-size} by quantifying the share of workers below, near, and above the hourly minimum-wage benchmark in 2025. Table~\ref{tab:minimum-wage-monthly-bunching-2025} reports the analogous exercise using implied monthly labor income relative to the monthly minimum wage. The alternative-reference tables use workers in firms with 101 or more workers or with 11--19 workers as the reference category. These changes alter the normalization but do not change the underlying empirical specifications. In the wage-level regressions, positive and negative size coefficients indicate higher and lower conditional wages relative to the omitted category. In the trend tables, the sign instead indicates whether that relative wage ratio rises or falls over time. In the interaction tables, the middle column is the difference between the formal and informal size coefficients, not a wage gap relative to the reference category on its own.

\begin{table}[htbp]
  \centering
  \caption{Distribution of hourly wages relative to the hourly minimum-wage benchmark, 2025}
  \label{tab:minimum-wage-bunching-2025}
  \small
  \begin{tabular}{lrrrrrrr}
    \toprule
    Firm size & Employment share & $<0.9$ & $0.9$--$1.1$ & $>1.1$ & P25 & Median & P75 \\
    \midrule
    \multicolumn{8}{l}{\textit{Formal workers}} \\
    Solo & 6.4\% & 22.2\% & 13.8\% & 64.0\% & 0.94 & 1.32 & 2.32 \\
    2--3 & 4.0\% & 21.3\% & 28.4\% & 50.3\% & 0.94 & 1.11 & 1.58 \\
    4--5 & 3.8\% & 14.2\% & 40.1\% & 45.7\% & 0.94 & 1.05 & 1.42 \\
    6--10 & 6.0\% & 11.6\% & 43.3\% & 45.1\% & 0.94 & 1.05 & 1.37 \\
    11--19 & 5.6\% & 9.0\% & 40.6\% & 50.5\% & 0.98 & 1.11 & 1.55 \\
    20--30 & 5.9\% & 7.3\% & 41.3\% & 51.4\% & 0.98 & 1.12 & 1.58 \\
    31--50 & 6.1\% & 7.5\% & 38.4\% & 54.0\% & 0.98 & 1.13 & 1.76 \\
    51--100 & 7.3\% & 6.2\% & 37.0\% & 56.8\% & 0.98 & 1.17 & 1.88 \\
    101+ & 54.9\% & 4.8\% & 28.1\% & 67.1\% & 1.02 & 1.38 & 2.53 \\
    \midrule
    \multicolumn{8}{l}{\textit{Informal workers}} \\
    Solo & 57.6\% & 74.3\% & 9.8\% & 15.9\% & 0.38 & 0.62 & 0.91 \\
    2--3 & 21.7\% & 67.0\% & 14.9\% & 18.2\% & 0.47 & 0.71 & 0.99 \\
    4--5 & 9.7\% & 61.1\% & 20.8\% & 18.1\% & 0.60 & 0.79 & 1.01 \\
    6--10 & 5.6\% & 53.6\% & 21.7\% & 24.6\% & 0.63 & 0.86 & 1.09 \\
    11--19 & 2.5\% & 49.3\% & 24.0\% & 26.6\% & 0.66 & 0.91 & 1.11 \\
    20--30 & 1.1\% & 48.3\% & 24.6\% & 27.0\% & 0.66 & 0.91 & 1.12 \\
    31--50 & 0.5\% & 47.6\% & 25.7\% & 26.7\% & 0.66 & 0.92 & 1.13 \\
    51--100 & 0.4\% & 42.3\% & 21.9\% & 35.8\% & 0.68 & 0.98 & 1.28 \\
    101+ & 0.9\% & 43.5\% & 19.5\% & 37.0\% & 0.66 & 0.95 & 1.37 \\
    \bottomrule
  \end{tabular}
  \vspace{0.3em}
  \begin{minipage}{0.95\textwidth}
  \footnotesize
  Notes: Wages are expressed relative to the hourly minimum-wage benchmark, computed as 1,423,500 pesos divided by $45\times52/12$ hours. This is an annual conversion benchmark, not a worker-specific legal compliance threshold. The columns $<0.9$, $0.9$--$1.1$, and $>1.1$ report the share within each size-formality group in each range. Employment shares are computed within each formality group. Quantiles refer to the wage-to-benchmark ratio.
  \end{minipage}
\end{table}


\begin{table}[htbp]
  \centering
  \caption{Distribution of monthly labor income relative to the monthly minimum wage, 2025}
  \label{tab:minimum-wage-monthly-bunching-2025}
  \small
  \begin{tabular}{lrrrrrrr}
    \toprule
    Firm size & Employment share & $<0.9$ & $0.9$--$1.1$ & $>1.1$ & P25 & Median & P75 \\
    \midrule
    \multicolumn{8}{l}{\textit{Formal workers}} \\
    Solo & 6.4\% & 20.1\% & 17.9\% & 62.0\% & 0.98 & 1.33 & 2.11 \\
    2--3 & 4.0\% & 9.6\% & 35.4\% & 55.0\% & 1.00 & 1.14 & 1.69 \\
    4--5 & 3.8\% & 3.4\% & 45.1\% & 51.5\% & 1.00 & 1.12 & 1.40 \\
    6--10 & 6.0\% & 2.7\% & 47.1\% & 50.2\% & 1.00 & 1.11 & 1.40 \\
    11--19 & 5.6\% & 1.8\% & 43.6\% & 54.5\% & 1.00 & 1.12 & 1.55 \\
    20--30 & 5.9\% & 1.4\% & 43.9\% & 54.7\% & 1.00 & 1.12 & 1.55 \\
    31--50 & 6.1\% & 1.4\% & 40.9\% & 57.7\% & 1.00 & 1.14 & 1.76 \\
    51--100 & 7.3\% & 1.1\% & 39.3\% & 59.6\% & 1.00 & 1.19 & 1.83 \\
    101+ & 54.9\% & 1.1\% & 29.4\% & 69.5\% & 1.00 & 1.40 & 2.46 \\
    \midrule
    \multicolumn{8}{l}{\textit{Informal workers}} \\
    Solo & 57.6\% & 78.0\% & 10.1\% & 11.9\% & 0.28 & 0.53 & 0.84 \\
    2--3 & 21.7\% & 65.5\% & 16.5\% & 18.0\% & 0.42 & 0.70 & 1.00 \\
    4--5 & 9.7\% & 56.3\% & 22.5\% & 21.3\% & 0.56 & 0.84 & 1.05 \\
    6--10 & 5.6\% & 50.9\% & 25.3\% & 23.7\% & 0.59 & 0.86 & 1.05 \\
    11--19 & 2.5\% & 46.2\% & 26.0\% & 27.7\% & 0.63 & 0.91 & 1.12 \\
    20--30 & 1.1\% & 46.3\% & 28.3\% & 25.4\% & 0.56 & 0.91 & 1.12 \\
    31--50 & 0.5\% & 42.3\% & 28.4\% & 29.3\% & 0.70 & 0.98 & 1.14 \\
    51--100 & 0.4\% & 41.9\% & 23.4\% & 34.7\% & 0.63 & 1.00 & 1.26 \\
    101+ & 0.9\% & 45.9\% & 24.9\% & 29.2\% & 0.56 & 0.95 & 1.26 \\
    \bottomrule
  \end{tabular}
  \vspace{0.3em}
  \begin{minipage}{0.95\textwidth}
  \footnotesize
  Notes: Monthly labor income refers to the main job. The columns $<0.9$, $0.9$--$1.1$, and $>1.1$ report the share within each size-formality group in each range relative to the statutory monthly minimum wage. Employment shares are computed within each formality group. Quantiles refer to the income-to-minimum ratio.
  \end{minipage}
\end{table}


\begin{table}[htbp]
  \centering
  \caption{Employment shares and mean real hourly income relative to 101+ workers, 2008 vs. 2025}
  \label{tab:descriptive-core-size-2008-2025-101-reference}
  \small
  \begin{tabular}{lrrrrrrr}
    \toprule
    & \multicolumn{2}{c}{Employment share} & \multicolumn{2}{c}{Mean hourly income} & \multicolumn{2}{c}{Wage ratio to 101+} & \multicolumn{1}{c}{Wage growth} \\
    \cmidrule(lr){2-3} \cmidrule(lr){4-5} \cmidrule(lr){6-7} \cmidrule(lr){8-8}
    Firm size & 2008 & 2025 & 2008 & 2025 & 2008 & 2025 & Annualized \\
    \midrule
    Solo & 36.1\% & 34.7\% &  5,388 &  6,534 & 0.42 & 0.39 & 1.1\% \\
    2--3 & 19.6\% & 13.8\% &  6,242 &  7,242 & 0.48 & 0.43 & 0.9\% \\
    4--5 & 7.3\% & 7.0\% &  7,196 &  8,242 & 0.56 & 0.49 & 0.8\% \\
    6--10 & 5.9\% & 5.8\% &  7,543 &  9,663 & 0.58 & 0.57 & 1.5\% \\
    11--19 & 3.4\% & 3.9\% &  8,679 & 11,634 & 0.67 & 0.69 & 1.7\% \\
    20--30 & 3.3\% & 3.3\% &  8,905 & 11,201 & 0.69 & 0.66 & 1.4\% \\
    31--50 & 2.6\% & 3.0\% & 10,424 & 12,574 & 0.81 & 0.75 & 1.1\% \\
    51--100 & 3.0\% & 3.5\% & 10,682 & 13,217 & 0.83 & 0.78 & 1.3\% \\
    101+ & 18.8\% & 25.2\% & 12,943 & 16,859 & 1.00 & 1.00 & 1.6\% \\
    \addlinespace
    All workers & 100.0\% & 100.0\% &  7,752 & 10,295 & 0.60 & 0.61 & 1.7\% \\
    \bottomrule
  \end{tabular}
  \vspace{0.3em}
  \begin{minipage}{0.95\textwidth}
  \footnotesize
  Notes: Mean hourly income is expressed in constant 2025 pesos. Wage ratios divide each firm-size mean by the mean among workers in firms with 101 or more workers in the same year. Annualized change is the compound annual change in mean hourly income between 2008 and 2025.
  \end{minipage}
\end{table}


\begin{table}[htbp]
  \centering
  \caption{Employment shares and mean real hourly income relative to 11--19 workers, 2008 vs. 2025}
  \label{tab:descriptive-core-size-2008-2025-11-19-reference}
  \small
  \begin{tabular}{lrrrrrrr}
    \toprule
    & \multicolumn{2}{c}{Employment share} & \multicolumn{2}{c}{Mean hourly income} & \multicolumn{2}{c}{Wage ratio to 11--19} & \multicolumn{1}{c}{Wage growth} \\
    \cmidrule(lr){2-3} \cmidrule(lr){4-5} \cmidrule(lr){6-7} \cmidrule(lr){8-8}
    Firm size & 2008 & 2025 & 2008 & 2025 & 2008 & 2025 & Annualized \\
    \midrule
    Solo & 36.1\% & 34.7\% &  5,388 &  6,534 & 0.62 & 0.56 & 1.1\% \\
    2--3 & 19.6\% & 13.8\% &  6,242 &  7,242 & 0.72 & 0.62 & 0.9\% \\
    4--5 & 7.3\% & 7.0\% &  7,196 &  8,242 & 0.83 & 0.71 & 0.8\% \\
    6--10 & 5.9\% & 5.8\% &  7,543 &  9,663 & 0.87 & 0.83 & 1.5\% \\
    11--19 & 3.4\% & 3.9\% &  8,679 & 11,634 & 1.00 & 1.00 & 1.7\% \\
    20--30 & 3.3\% & 3.3\% &  8,905 & 11,201 & 1.03 & 0.96 & 1.4\% \\
    31--50 & 2.6\% & 3.0\% & 10,424 & 12,574 & 1.20 & 1.08 & 1.1\% \\
    51--100 & 3.0\% & 3.5\% & 10,682 & 13,217 & 1.23 & 1.14 & 1.3\% \\
    101+ & 18.8\% & 25.2\% & 12,943 & 16,859 & 1.49 & 1.45 & 1.6\% \\
    \addlinespace
    All workers & 100.0\% & 100.0\% &  7,752 & 10,295 & 0.89 & 0.88 & 1.7\% \\
    \bottomrule
  \end{tabular}
  \vspace{0.3em}
  \begin{minipage}{0.95\textwidth}
  \footnotesize
  Notes: Mean hourly income is expressed in constant 2025 pesos. Wage ratios divide each firm-size mean by the mean among workers in firms with 11--19 workers in the same year. Annualized change is the compound annual change in mean hourly income between 2008 and 2025.
  \end{minipage}
\end{table}


\begin{table}[htbp]
  \centering
  \caption{Firm-size wage premium regressions: 101+ workers as the reference category}
  \label{tab:firm-size-wage-premium-regressions-101-reference}
  \small
  \begin{tabular}{lcccc}
    \toprule
    & \multicolumn{4}{c}{Dependent variable: log real hourly labor income} \\
    \cmidrule(lr){2-5}
    & (1) & (2) & (3) & (4) \\
    \midrule
    Firm size: Solo & -1.028*** & -0.832*** & -0.613*** & -0.372*** \\
     & (0.102) & (0.068) & (0.062) & (0.060) \\
    Firm size: 2-3 & -0.838*** & -0.584*** & -0.398*** & -0.179*** \\
     & (0.095) & (0.056) & (0.046) & (0.042) \\
    Firm size: 4-5 & -0.656*** & -0.414*** & -0.244*** & -0.056* \\
     & (0.098) & (0.052) & (0.040) & (0.032) \\
    Firm size: 6-10 & -0.522*** & -0.323*** & -0.184*** & -0.049 \\
     & (0.092) & (0.051) & (0.037) & (0.032) \\
    Firm size: 11-19 & -0.394*** & -0.247*** & -0.147*** & -0.062 \\
     & (0.081) & (0.054) & (0.041) & (0.038) \\
    Firm size: 20-30 & -0.330*** & -0.217*** & -0.139*** & -0.090** \\
     & (0.069) & (0.050) & (0.038) & (0.038) \\
    Firm size: 31-50 & -0.258*** & -0.169*** & -0.113*** & -0.090** \\
     & (0.059) & (0.045) & (0.033) & (0.034) \\
    Firm size: 51-100 & -0.199*** & -0.121*** & -0.082*** & -0.072** \\
     & (0.051) & (0.036) & (0.027) & (0.028) \\
    \midrule
    Gender, age, and education controls & No & No & Yes & Yes \\
    Formality control & No & No & No & Yes \\
    Sector-year fixed effects & No & Yes & Yes & Yes \\
    Observations & 5,426,316 & 5,426,316 & 5,426,316 & 5,426,316 \\
    $R^2$ & 0.218 & 0.277 & 0.381 & 0.391 \\
    \bottomrule
  \end{tabular}
  \vspace{0.3em}
  \begin{minipage}{0.95\textwidth}
  \footnotesize
  Notes: Specifications as in Table~\ref{tab:firm-size-wage-premium-regressions}, with 101+ workers as the omitted category. Parentheses contain sector-clustered standard errors. * $p<0.10$, ** $p<0.05$, *** $p<0.01$.
  \end{minipage}
\end{table}


\begin{table}[htbp]
  \centering
  \caption{Firm-size wage premium regressions: 11--19 workers as the reference category}
  \label{tab:firm-size-wage-premium-regressions-11-19-reference}
  \small
  \begin{tabular}{lcccc}
    \toprule
    & \multicolumn{4}{c}{Dependent variable: log real hourly labor income} \\
    \cmidrule(lr){2-5}
    & (1) & (2) & (3) & (4) \\
    \midrule
    Firm size: Solo & -0.633*** & -0.585*** & -0.466*** & -0.310*** \\
     & (0.035) & (0.041) & (0.042) & (0.051) \\
    Firm size: 2-3 & -0.444*** & -0.337*** & -0.252*** & -0.117*** \\
     & (0.029) & (0.033) & (0.028) & (0.030) \\
    Firm size: 4-5 & -0.262*** & -0.167*** & -0.097*** & 0.006 \\
     & (0.029) & (0.021) & (0.021) & (0.020) \\
    Firm size: 6-10 & -0.128*** & -0.076*** & -0.037*** & 0.013 \\
     & (0.017) & (0.010) & (0.010) & (0.010) \\
    Firm size: 20-30 & 0.064*** & 0.030*** & 0.008 & -0.028*** \\
     & (0.015) & (0.009) & (0.006) & (0.005) \\
    Firm size: 31-50 & 0.136*** & 0.078*** & 0.034*** & -0.028*** \\
     & (0.026) & (0.015) & (0.010) & (0.008) \\
    Firm size: 51-100 & 0.196*** & 0.126*** & 0.065*** & -0.010 \\
     & (0.034) & (0.023) & (0.017) & (0.014) \\
    Firm size: 101+ & 0.394*** & 0.247*** & 0.147*** & 0.062 \\
     & (0.081) & (0.054) & (0.041) & (0.038) \\
    \midrule
    Gender, age, and education controls & No & No & Yes & Yes \\
    Formality control & No & No & No & Yes \\
    Sector-year fixed effects & No & Yes & Yes & Yes \\
    Observations & 5,426,316 & 5,426,316 & 5,426,316 & 5,426,316 \\
    $R^2$ & 0.218 & 0.277 & 0.381 & 0.391 \\
    \bottomrule
  \end{tabular}
  \vspace{0.3em}
  \begin{minipage}{0.95\textwidth}
  \footnotesize
  Notes: Specifications as in Table~\ref{tab:firm-size-wage-premium-regressions}, with 11--19 workers as the omitted category. Parentheses contain sector-clustered standard errors. * $p<0.10$, ** $p<0.05$, *** $p<0.01$.
  \end{minipage}
\end{table}


\begin{table}[htbp]
  \centering
  \caption{Testing whether wage gaps relative to 101+ workers changed over time}
  \label{tab:firm-size-premium-trend-test-101-reference}
  \small
  \begin{tabular}{lcccc}
    \toprule
    Firm size & Annual trend & S.E. & $p$-value & Implied 2008--2025 change \\
    \midrule
    Solo & -0.0045 & (0.0043) & 0.313 & -7.3\% \\
    2-3 & -0.0012 & (0.0023) & 0.612 & -2.0\% \\
    4-5 & -0.0013 & (0.0023) & 0.587 & -2.1\% \\
    6-10 & -0.0023 & (0.0017) & 0.190 & -3.9\% \\
    11-19 & -0.0017 & (0.0012) & 0.175 & -2.9\% \\
    20-30 & -0.0026 & (0.0009) & 0.012 & -4.3\% \\
    31-50 & -0.0017 & (0.0009) & 0.065 & -2.9\% \\
    51-100 & -0.0018 & (0.0013) & 0.181 & -3.1\% \\
    \bottomrule
  \end{tabular}
  \vspace{0.3em}
  \begin{minipage}{0.95\textwidth}
  \footnotesize
  Notes: Annual changes in log-wage gaps from Equation~(\ref{eq:trend-benchmark}), expressed relative to 101+ workers. The final column is the compounded 2008--2025 percentage change in that wage ratio, not a percentage-point change in the premium. Two-sided tests of zero trends use sector-clustered standard errors.
  \end{minipage}
\end{table}


\begin{table}[htbp]
  \centering
  \caption{Testing whether wage gaps relative to 11--19 workers changed over time}
  \label{tab:firm-size-premium-trend-test-11-19-reference}
  \small
  \begin{tabular}{lcccc}
    \toprule
    Firm size & Annual trend & S.E. & $p$-value & Implied 2008--2025 change \\
    \midrule
    Solo & -0.0027 & (0.0039) & 0.491 & -4.5\% \\
    2-3 & 0.0005 & (0.0016) & 0.746 & 0.9\% \\
    4-5 & 0.0005 & (0.0014) & 0.736 & 0.8\% \\
    6-10 & -0.0006 & (0.0011) & 0.611 & -1.0\% \\
    20-30 & -0.0009 & (0.0010) & 0.392 & -1.5\% \\
    31-50 & 0.0000 & (0.0012) & 0.988 & 0.0\% \\
    51-100 & -0.0001 & (0.0014) & 0.940 & -0.2\% \\
    101+ & 0.0017 & (0.0012) & 0.175 & 3.0\% \\
    \bottomrule
  \end{tabular}
  \vspace{0.3em}
  \begin{minipage}{0.95\textwidth}
  \footnotesize
  Notes: Annual changes in log-wage gaps from Equation~(\ref{eq:trend-benchmark}), expressed relative to 11--19 workers. The final column is the compounded 2008--2025 percentage change in that wage ratio, not a percentage-point change in the premium. Two-sided tests of zero trends use sector-clustered standard errors.
  \end{minipage}
\end{table}


\begin{table}[htbp]
  \centering
  \caption{Firm-size coefficients by formality status: 101+ workers as reference}
  \label{tab:formality-interaction-coefficients-101-reference}
  \small
  \begin{tabular}{lccc}
    \toprule
    Firm size & $\hat{\beta}_g$ & $\hat{\theta}_g$ & $\hat{\beta}_g+\hat{\theta}_g$ \\
    \midrule
    Solo & -0.336*** & 0.182** & -0.154 \\
     & (0.046) & (0.083) & (0.088) \\
    2-3 & -0.139*** & 0.007 & -0.132*** \\
     & (0.034) & (0.047) & (0.040) \\
    4-5 & 0.007 & -0.136** & -0.129*** \\
     & (0.031) & (0.048) & (0.033) \\
    6-10 & 0.056* & -0.187*** & -0.131*** \\
     & (0.029) & (0.046) & (0.035) \\
    11-19 & 0.072** & -0.191*** & -0.119** \\
     & (0.027) & (0.044) & (0.041) \\
    20-30 & 0.066** & -0.196*** & -0.130*** \\
     & (0.024) & (0.042) & (0.040) \\
    31-50 & 0.064** & -0.175*** & -0.112*** \\
     & (0.022) & (0.041) & (0.036) \\
    51-100 & 0.073*** & -0.158*** & -0.085** \\
     & (0.021) & (0.039) & (0.030) \\
    \bottomrule
  \end{tabular}
  \vspace{0.3em}
  \begin{minipage}{0.92\textwidth}
  \footnotesize
  Notes: Log-point coefficients from Equation~(\ref{eq:formality-interaction}), with 101+ workers as the reference within each formality group and tildes suppressed. The first and third columns are informal and formal size gaps. The middle column is the formal-minus-informal difference in those size gaps. Parentheses contain sector-clustered standard errors. * $p<0.10$, ** $p<0.05$, *** $p<0.01$.
  \end{minipage}
\end{table}


\begin{table}[htbp]
  \centering
  \caption{Firm-size coefficients by formality status: 11--19 workers as reference}
  \label{tab:formality-interaction-coefficients-11-19-reference}
  \small
  \begin{tabular}{lccc}
    \toprule
    Firm size & $\hat{\beta}_g$ & $\hat{\theta}_g$ & $\hat{\beta}_g+\hat{\theta}_g$ \\
    \midrule
    Solo & -0.415*** & 0.372*** & -0.043 \\
     & (0.038) & (0.075) & (0.079) \\
    2-3 & -0.212*** & 0.191*** & -0.021 \\
     & (0.025) & (0.022) & (0.024) \\
    4-5 & -0.063*** & 0.049** & -0.013 \\
     & (0.019) & (0.018) & (0.018) \\
    6-10 & -0.016 & 0.002 & -0.014 \\
     & (0.010) & (0.012) & (0.012) \\
    20-30 & -0.005 & -0.004 & -0.009* \\
     & (0.007) & (0.009) & (0.005) \\
    31-50 & -0.009 & 0.019 & 0.010 \\
     & (0.009) & (0.012) & (0.008) \\
    51-100 & 0.004 & 0.032 & 0.036** \\
     & (0.023) & (0.026) & (0.015) \\
    101+ & -0.074** & 0.198*** & 0.124*** \\
     & (0.027) & (0.044) & (0.040) \\
    \bottomrule
  \end{tabular}
  \vspace{0.3em}
  \begin{minipage}{0.92\textwidth}
  \footnotesize
  Notes: Log-point coefficients from Equation~(\ref{eq:formality-interaction}), with 11--19 workers as the reference within each formality group and tildes suppressed. The first and third columns are informal and formal size gaps. The middle column is the formal-minus-informal difference in those size gaps. Parentheses contain sector-clustered standard errors. * $p<0.10$, ** $p<0.05$, *** $p<0.01$.
  \end{minipage}
\end{table}


Table~\ref{tab:firm-size-local-geography-robustness} adds geographic and occupational-position controls to the common-premium comparison model. Its final column absorbs sector-department-year fixed effects and controls for occupational position. It constrains the size premium to be the same for formal and informal workers. Under this restriction, workers in firms with 101 or more workers earn about 40.2 percent more than observably comparable solo workers, a statistically significant difference. Tables~\ref{tab:firm-size-local-geography-robustness-101-reference} and \ref{tab:firm-size-local-geography-robustness-11-19-reference} report the same specifications with alternative reference categories. The preferred model relaxes the common-premium restriction and is reported in Table~\ref{tab:firm-size-local-geography-formality-robustness} in the main text.

\begin{table}[htbp]
  \centering
  \caption{Common-premium estimates with geographic and occupational controls}
  \label{tab:firm-size-local-geography-robustness}
  \small
  \begin{tabular}{lcccc}
    \toprule
    & \multicolumn{4}{c}{Dependent variable: log real hourly labor income} \\
    \cmidrule(lr){2-5}
    & (1) & (2) & (3) & (4) \\
    \midrule
    Firm size: 2-3 & 0.188*** & 0.168*** & 0.167*** & 0.102*** \\
     & (0.011) & (0.011) & (0.011) & (0.013) \\
    Firm size: 4-5 & 0.307*** & 0.277*** & 0.273*** & 0.225*** \\
     & (0.014) & (0.014) & (0.014) & (0.015) \\
    Firm size: 6-10 & 0.316*** & 0.282*** & 0.278*** & 0.265*** \\
     & (0.018) & (0.017) & (0.017) & (0.016) \\
    Firm size: 11-19 & 0.302*** & 0.272*** & 0.267*** & 0.272*** \\
     & (0.018) & (0.017) & (0.017) & (0.015) \\
    Firm size: 20-30 & 0.271*** & 0.241*** & 0.237*** & 0.258*** \\
     & (0.019) & (0.017) & (0.017) & (0.015) \\
    Firm size: 31-50 & 0.271*** & 0.245*** & 0.240*** & 0.263*** \\
     & (0.018) & (0.017) & (0.017) & (0.015) \\
    Firm size: 51-100 & 0.289*** & 0.262*** & 0.257*** & 0.282*** \\
     & (0.018) & (0.017) & (0.017) & (0.015) \\
    Firm size: 101+ & 0.359*** & 0.338*** & 0.325*** & 0.338*** \\
     & (0.023) & (0.024) & (0.024) & (0.022) \\
    \midrule
    Worker and formality controls & Yes & Yes & Yes & Yes \\
    Occupational-position controls & No & No & No & Yes \\
    Fixed effects & $s\times t$ & $s\times t$, $d\times t$ & $s\times d\times t$ & $s\times d\times t$ \\
    Observations & 5,268,514 & 5,268,514 & 5,268,375 & 5,268,375 \\
    $R^2$ & 0.386 & 0.404 & 0.414 & 0.429 \\
    \bottomrule
  \end{tabular}
  \vspace{0.3em}
  \begin{minipage}{0.95\textwidth}
  \footnotesize
  Notes: The omitted category is solo workers. Column (4) corresponds to Equation~(\ref{eq:common-premium-geography}). $s$, $d$, and $t$ denote sector, department, and year. Parentheses contain sector-department-clustered standard errors. * $p<0.10$, ** $p<0.05$, *** $p<0.01$.
  \end{minipage}
\end{table}


\begin{table}[htbp]
  \centering
  \caption{Common-premium estimates with geographic and status-in-employment controls: 101+ workers as reference}
  \label{tab:firm-size-local-geography-robustness-101-reference}
  \small
  \begin{tabular}{lcccc}
    \toprule
    & \multicolumn{4}{c}{Dependent variable: log real hourly labor income} \\
    \cmidrule(lr){2-5}
    & (1) & (2) & (3) & (4) \\
    \midrule
    Firm size: Solo & -0.372*** & -0.352*** & -0.337*** & -0.343*** \\
     & (0.024) & (0.024) & (0.024) & (0.022) \\
    Firm size: 2-3 & -0.179*** & -0.178*** & -0.165*** & -0.242*** \\
     & (0.018) & (0.018) & (0.017) & (0.016) \\
    Firm size: 4-5 & -0.056*** & -0.066*** & -0.056*** & -0.117*** \\
     & (0.013) & (0.014) & (0.014) & (0.014) \\
    Firm size: 6-10 & -0.049*** & -0.063*** & -0.053*** & -0.078*** \\
     & (0.012) & (0.013) & (0.013) & (0.012) \\
    Firm size: 11-19 & -0.062*** & -0.072*** & -0.064*** & -0.071*** \\
     & (0.014) & (0.015) & (0.015) & (0.013) \\
    Firm size: 20-30 & -0.090*** & -0.100*** & -0.091*** & -0.082*** \\
     & (0.014) & (0.015) & (0.014) & (0.012) \\
    Firm size: 31-50 & -0.090*** & -0.096*** & -0.088*** & -0.076*** \\
     & (0.012) & (0.014) & (0.013) & (0.011) \\
    Firm size: 51-100 & -0.072*** & -0.079*** & -0.070*** & -0.058*** \\
     & (0.012) & (0.013) & (0.012) & (0.011) \\
    \midrule
    Worker and formality controls & Yes & Yes & Yes & Yes \\
    Status-in-employment controls & No & No & No & Yes \\
    Fixed effects & $s\times t$ & $s\times t$, $d\times t$ & $s\times d\times t$ & $s\times d\times t$ \\
    Observations & 5,426,316 & 5,426,316 & 5,426,174 & 5,426,174 \\
    $R^2$ & 0.391 & 0.410 & 0.420 & 0.435 \\
    \bottomrule
  \end{tabular}
  \vspace{0.3em}
  \begin{minipage}{0.95\textwidth}
  \footnotesize
  Notes: Specifications as in Table~\ref{tab:firm-size-local-geography-robustness}, with 101+ workers as the omitted category. $s$, $d$, and $t$ denote sector, department, and year. Parentheses contain sector-department-clustered standard errors. * $p<0.10$, ** $p<0.05$, *** $p<0.01$.
  \end{minipage}
\end{table}


\begin{table}[htbp]
  \centering
  \caption{Common-premium estimates with geographic and occupational controls: 11--19 workers as reference}
  \label{tab:firm-size-local-geography-robustness-11-19-reference}
  \small
  \begin{tabular}{lcccc}
    \toprule
    & \multicolumn{4}{c}{Dependent variable: log real hourly labor income} \\
    \cmidrule(lr){2-5}
    & (1) & (2) & (3) & (4) \\
    \midrule
    Firm size: Solo & -0.302*** & -0.272*** & -0.267*** & -0.272*** \\
     & (0.018) & (0.017) & (0.017) & (0.015) \\
    Firm size: 2-3 & -0.114*** & -0.104*** & -0.100*** & -0.169*** \\
     & (0.013) & (0.011) & (0.011) & (0.009) \\
    Firm size: 4-5 & 0.005 & 0.005 & 0.006 & -0.047*** \\
     & (0.008) & (0.008) & (0.008) & (0.007) \\
    Firm size: 6-10 & 0.014*** & 0.010** & 0.011** & -0.007 \\
     & (0.005) & (0.005) & (0.005) & (0.004) \\
    Firm size: 20-30 & -0.031*** & -0.031*** & -0.030*** & -0.014*** \\
     & (0.004) & (0.004) & (0.004) & (0.004) \\
    Firm size: 31-50 & -0.031*** & -0.027*** & -0.027*** & -0.008* \\
     & (0.005) & (0.005) & (0.005) & (0.004) \\
    Firm size: 51-100 & -0.014* & -0.010 & -0.010 & 0.010 \\
     & (0.007) & (0.007) & (0.007) & (0.007) \\
    Firm size: 101+ & 0.057*** & 0.066*** & 0.058*** & 0.066*** \\
     & (0.014) & (0.015) & (0.015) & (0.013) \\
    \midrule
    Worker and formality controls & Yes & Yes & Yes & Yes \\
    Occupational-position controls & No & No & No & Yes \\
    Fixed effects & $s\times t$ & $s\times t$, $d\times t$ & $s\times d\times t$ & $s\times d\times t$ \\
    Observations & 5,268,514 & 5,268,514 & 5,268,375 & 5,268,375 \\
    $R^2$ & 0.386 & 0.404 & 0.414 & 0.429 \\
    \bottomrule
  \end{tabular}
  \vspace{0.3em}
  \begin{minipage}{0.95\textwidth}
  \footnotesize
  Notes: Specifications as in Table~\ref{tab:firm-size-local-geography-robustness}, with 11--19 workers as the omitted category. $s$, $d$, and $t$ denote sector, department, and year. Parentheses contain sector-department-clustered standard errors. * $p<0.10$, ** $p<0.05$, *** $p<0.01$.
  \end{minipage}
\end{table}


\FloatBarrier
Table~\ref{tab:income-counterfactual-sensitivity} extends the income scenarios
in Section~\ref{subsec:income-counterfactuals}. Preserving observed log-wage
residuals gives aggregate changes close to those based on fitted income.
For example, the combined formal-101+ scenario yields 25.0 percent rather than
25.4 percent. The last two rows select the highest estimated size component
within each formality group. The selected destinations are 101+ for formal
workers and 51--100 for informal workers. Assigning informal workers to
51--100 raises total predicted income by 10.0 percent. Assigning both groups
to their respective selected destinations raises it by 13.7 percent without
changing formality. This exercise includes movement from informal 101+ to
51--100 and therefore is not uniformly a move toward larger firms. The
destinations maximize the model's fitted income within each formality group.
Because the size coefficients do not interact with other worker characteristics,
the selected size is the same for all workers of a given formality status,
not an individually estimated best job.

% Generated by replication/build_counterfactual_tables.py.
\begin{table}[p]
  \centering
  \caption{Income scenarios with alternative baselines and destinations}
  \label{tab:income-counterfactual-sensitivity}
  \small
  \begin{tabular}{>{\raggedright\arraybackslash}p{0.52\textwidth}cc}
    \toprule
    Scenario & Predicted baseline & Observed baseline \\
    & (\%) & (\%) \\
    \midrule
    Informal to formal, same size & 13.5 & 13.3 \\[0.4em]
    Formal workers to formal 101+ & 3.8 & 3.8 \\[0.4em]
    Informal workers to informal 101+ & 7.6 & 7.5 \\[0.4em]
    All workers to formal 101+ & 25.9 & 25.5 \\[0.4em]
    \midrule
    \multicolumn{3}{l}{\textit{Destinations selected from the estimated profiles}} \\
    Informal workers to informal 51--100 & 10.1 & 9.9 \\[0.4em]
    Formal to 101+, informal to 51--100 & 13.8 & 13.7 \\[0.4em]
    \bottomrule
  \end{tabular}
  \vspace{0.4em}
  \begin{minipage}{0.97\textwidth}
  \footnotesize
  Notes: Percentage changes in total weekly labor income for the entire 2025 sample. The predicted baseline uses Equation~(\ref{eq:income-counterfactual}). The observed baseline rescales observed income by the same estimated wage ratios, preserving log residuals. Both use pooled coefficients from Equation~(\ref{eq:main}). In the last two rows, destinations maximize the estimated size component within each formality group, including moves from informal 101+ to informal 51--100.
  \end{minipage}
\end{table}


Table~\ref{tab:income-counterfactual-support} reports the share of intended
movers with no observed destination counterpart and scenarios restricted by
the number of destination observations. For each intended move, we
check whether workers with the same sector, department, and occupational
position are observed in the destination size-formality category. Panel A
shows that absent destinations are especially common for assignments to
large informal categories. Panel B leaves workers at their origin whenever
the destination has fewer than the stated number of observations. Requiring
30 observations in the pooled sample reduces the total predicted income
change for assignment to informal 101+ from 7.7 to 3.7 percent. Requiring
30 observations in 2025 permits no such moves. There are 1,843 informal 101+
observations in 2025, but no destination sector-department-occupation cell
contains more than 25. The restricted scenarios thus change who is assigned
to a destination, not the estimated coefficients. Their smaller changes do
not establish smaller returns for an otherwise identical reallocation.
The screens are diagnostics of extrapolation, not tests that a job transition
is feasible or that workers are comparable on all characteristics.

% Generated by replication/build_counterfactual_tables.py.
\begin{table}[p]
  \centering
  \caption{Observed destination groups and sensitivity to extrapolation, 2025}
  \label{tab:income-counterfactual-support}
  \small
  \begin{tabular}{>{\raggedright\arraybackslash}p{0.40\textwidth}cccc}
    \toprule
    \multicolumn{5}{l}{\textit{A. Intended moves and absent destination cells (\%)}} \\
    \midrule
    Scenario & \shortstack{Intended\\movers} & \multicolumn{2}{c}{\shortstack{Absent destination\\in 2025}} & \shortstack{Absent\\in pooled} \\
    \midrule
    Informal to formal, same size & 55.1 & \multicolumn{2}{c}{6.0} & 0.4 \\[0.3em]
    Formal to 101+ & 20.2 & \multicolumn{2}{c}{6.9} & 1.8 \\[0.3em]
    Informal to 101+ & 54.6 & \multicolumn{2}{c}{39.0} & 6.9 \\[0.3em]
    All to formal 101+ & 75.3 & \multicolumn{2}{c}{17.2} & 4.6 \\[0.3em]
    Informal to 51--100 & 54.9 & \multicolumn{2}{c}{55.0} & 10.8 \\[0.3em]
    Formal to 101+, informal to 51--100 & 75.1 & \multicolumn{2}{c}{42.0} & 8.4 \\[0.3em]
    \midrule
    \multicolumn{5}{l}{\textit{B. Change in total predicted income with destination screens (\%)}} \\
    \midrule
    Scenario & No screen & \shortstack{Pooled\\$n\geq30$} & \shortstack{2025\\$n\geq10$} & \shortstack{2025\\$n\geq30$} \\
    \midrule
    Informal to formal, same size & 13.5 & 12.0 & 8.7 & 5.2\\[0.3em]
    Formal to 101+ & 3.8 & 3.2 & 2.9 & 2.4\\[0.3em]
    Informal to 101+ & 7.6 & 3.8 & 0.5 & \textit{None}\\[0.3em]
    All to formal 101+ & 25.9 & 17.7 & 11.4 & 6.7\\[0.3em]
    Informal to 51--100 & 10.1 & 2.3 & $<0.1$ & \textit{None}\\[0.3em]
    Formal to 101+, informal to 51--100 & 13.8 & 5.4 & 2.9 & 2.4\\[0.3em]
    \bottomrule
  \end{tabular}
  \vspace{0.4em}
  \begin{minipage}{0.97\textwidth}
  \footnotesize
  Notes: Destination cells are defined by sector, department, status in employment, size, and formality. Panel A reports intended movers as a percentage of all workers. The absent-destination columns instead use intended movers as their denominator. Pooled counts use all estimation years. Panel B permits a move only when its destination meets the stated unweighted observation count. Otherwise the worker stays at origin. The all-worker income denominator is unchanged. \textit{None} means that no intended move qualifies. For the joint selected-destination scenario, the 2025 $n\geq30$ column contains only formal moves.
  \end{minipage}
\end{table}

\FloatBarrier
